\textbf{Abstract:} Graph neural networks assume every node carries an informative feature vector, yet in practice many nodes have no observed features at all, such as users without profiles, newly listed items, or failed sensors, and message passing degrades sharply when neighborhoods are empty. Prior work largely treats this as feature imputation, reconstructing missing attributes before classification, and adds self-supervision by analogy without studying which auxiliary objective actually helps. We introduce NESS, a graph neural network for learning under missing node features that couples a Feature-Propagation prefill, giving every missing node long-range structural signal, with a set of pluggable self-supervised objectives on a message-passing encoder. Evaluating by downstream node-classification accuracy rather than reconstruction quality, NESS outperforms imputation-based, imputation-free, and per-node-parameter baselines, with the largest advantage under severe missingness and at a scale where generative imputation is infeasible. Beyond the method, we conduct an extensive study of self-supervised objectives for learning under missingness. We find that self-supervision generally helps, but the size of its benefit depends on the setting: it contributes most when the encoder is information-starved under severe missingness, and its benefit varies systematically with the characteristics of the objective. The label-aware neighbor-histogram is the most reliable objective, and the effective objective transfers as an enhancement to other backbones. These findings offer a principled basis for designing self-supervision under missing node features.
